\chapter{МЕТОДОЛОГИЯ СРАВНИТЕЛЬНОГО АНАЛИЗА}
\label{ch:methodology}

\section{Математические модели тестовых систем}
\label{sec:systems}

Для проведения объективного сравнительного анализа методов DMD, SINDy и EDMD необходимо сформировать унифицированные наборы синтетических данных. Использование синтетических данных позволяет точно контролировать наличие шума, объем выборки и параметры системы, что необходимо для верификации алгоритмов. В качестве бенчмарков выбраны три системы с различным характером динамики: линейная система с управлением (LTI), нелинейный осциллятор Ван-дер-Поля с управлением и хаотическая система Лоренца.

\subsection{Линейная стационарная система (LTI) с управлением}

Линейная стационарная система представляет собой базовый тест для проверки работоспособности алгоритмов. В дискретном времени динамика системы описывается уравнением:
\begin{equation}
\label{eq:lti}
x_{k+1} = A x_k + B u_k,
\end{equation}
где $x_k \in \mathbb{R}^n$ --- вектор состояния системы в момент времени $k$, $u_k \in \mathbb{R}^m$ --- вектор управляющего воздействия, $A \in \mathbb{R}^{n \times n}$ --- матрица собственной динамики, $B \in \mathbb{R}^{n \times m}$ --- матрица управления. 

\subsection{Осциллятор Ван-дер-Поля с управлением}

Осциллятор Ван-дер-Поля является классическим примером нелинейной системы с затуханием, зависящим от координаты. Модель с внешним управлением описывается дифференциальным уравнением второго порядка:
\begin{equation}
\label{eq:vdp_second_order}
\ddot{x} - \mu(1-x^2)\dot{x} + x = u(t),
\end{equation}
где $\mu > 0$ --- параметр нелинейного затухания, определяющий жесткость предельного цикла, $u(t)$ --- внешнее управляющее воздействие. 

Для применения алгоритмов восстановления динамики уравнение \ref{eq:vdp_second_order} приводится к системе дифференциальных уравнений первого порядка путем введения переменных состояния $x_1 = x$ и $x_2 = \dot{x}$:
\begin{equation}
\label{eq:vdp_state_space}
\begin{cases}
\dot{x}_1 = x_2 \\
\dot{x}_2 = \mu(1-x_1^2)x_2 - x_1 + u(t)
\end{cases}.
\end{equation}

\subsection{Система Лоренца}

Для тестирования алгоритмов на системах с хаотической динамикой и странными аттракторами используется система Лоренца. Она описывается системой трех обыкновенных нелинейных дифференциальных уравнений:
\begin{equation}
\label{eq:lorenz}
\begin{cases}
\dot{x} = \sigma(y - x) \\
\dot{y} = x(\rho - z) - y \\
\dot{z} = xy - \beta z
\end{cases},
\end{equation}
где $x, y, z$ --- переменные состояния, а параметры системы фиксируются на классических значениях $\sigma = 10$, $\rho = 28$, $\beta = 8/3$, при которых наблюдается детерминированный хаос.

\subsection{Модель наблюдений и добавление шума}

В реальных задачах данные всегда зашумлены. Чтобы приблизить синтетические данные к реальным условиям эксплуатации, вводится аддитивный измерительный шум. Наблюдаемые векторы состояния $x^{obs}_k$ формируются по следующему правилу:
\begin{equation}
\label{eq:noise_model}
x^{obs}_k = x(t_k) + \epsilon_k, \quad \epsilon_k \sim \mathcal{N}(0, \alpha^2 I),
\end{equation}
где $x(t_k)$ --- истинное значение вектора состояния, полученное путем численного интегрирования систем уравнений на сетке времени $t_k$, $\epsilon_k$ --- вектор гауссовского белого шума с нулевым математическим ожиданием и ковариационной матрицей $\alpha^2 I$, где $\alpha$ --- амплитуда (стандартное отклонение) шума.